% Zumdahl Chapter 12 test -- 20 MCQ + 2 FRQ. 50 min. All sections on-CED.
\documentclass{apchem-exam}

\apcunitword{Chapter}
\apcunit{12}
\apcunittitle{Chemical Kinetics}
\apcdoctitle{Chapter Test}
\apcsubtitle{Zumdahl \S12.1--12.7 \textbullet\ 50 minutes \textbullet\ kinetics equations provided with Section II}

\begin{document}
\apctitleblock

\examsection{I}{Multiple Choice}{20 questions \textbullet\ 30 minutes \textbullet\ 1 point each}{\nocalculator}

% ---- rates -----------------------------------------------------------------
\begin{mcq}
  For \ce{2 A -> 3 B}, if A disappears at \SI{0.10}{\Molar\per\second}, B
  appears at
  \choices
    \choice{\SI{0.10}{\Molar\per\second}}
    \correct{\SI{0.15}{\Molar\per\second}}
    \choice{\SI{0.067}{\Molar\per\second}}
    \choice{\SI{0.30}{\Molar\per\second}}
\end{mcq}
\why{$\tfrac32 \times 0.10$ --- three B for every two A.}
\distractor{C}{inverted the ratio}

\begin{mcq}
  The instantaneous rate of a reaction at a given moment is
  \choices
    \choice{the slope of the line joining two points on the curve}
    \correct{the slope of the tangent to the concentration--time curve}
    \choice{always constant}
    \choice{the same as the initial rate}
\end{mcq}
\why{Average rate uses a chord; instantaneous uses a tangent.}

\begin{mcq}
  Initial rates are used to determine rate laws because at $t = 0$
  \choices
    \correct{concentrations are known exactly and no products have
      accumulated}
    \choice{the rate is at its slowest}
    \choice{the rate constant is largest}
    \choice{temperature is lowest}
\end{mcq}
\why{Accumulated product could drive a reverse reaction and confound the
measurement.}

% ---- rate laws -------------------------------------------------------------
\begin{mcq}
  The exponents in a rate law are determined by
  \choices
    \choice{the coefficients in the balanced equation}
    \correct{experiment}
    \choice{the molar masses of the reactants}
    \choice{the sign of $\Delta H$}
\end{mcq}
\why{The central rule of the chapter --- orders reflect the mechanism, not
the stoichiometry.}

\begin{mcq}
  For rate $= k[\ce{A}]^2[\ce{B}]$, doubling both concentrations changes the
  rate by a factor of
  \choices
    \choice{4}
    \choice{6}
    \correct{8}
    \choice{2}
\end{mcq}
\why{$2^2 \times 2^1 = 8$.}

\begin{mcq}
  If doubling the concentration of a reactant leaves the rate unchanged, the
  order with respect to that reactant is
  \choices
    \correct{0}
    \choice{1}
    \choice{2}
    \choice{$-1$}
\end{mcq}
\why{Zero order means the rate is independent of that concentration.}

\begin{mcq}
  A reaction is first order in A and second order in B. The overall order is
  \choices
    \choice{1}
    \choice{2}
    \correct{3}
    \choice{4}
\end{mcq}
\why{Orders add.}

\begin{mcq}
  \ce{2 N2O5 -> 4 NO2 + O2} is found to be first order in \ce{N2O5}. This
  shows that
  \choices
    \choice{the balanced equation was written incorrectly}
    \correct{reaction order cannot be predicted from coefficients}
    \choice{the reaction has no mechanism}
    \choice{the rate is constant}
\end{mcq}
\why{Coefficient 2, order 1 --- exactly the point.}

% ---- integrated rate laws --------------------------------------------------
\begin{mcq}
  A plot of $\ln[\ce{A}]$ versus time gives a straight line. The reaction is
  \choices
    \choice{zero order}
    \correct{first order}
    \choice{second order}
    \choice{third order}
\end{mcq}
\why{Each order has exactly one linear plot.}

\begin{mcq}
  A plot of $1/[\ce{A}]$ versus time gives a straight line. The slope equals
  \choices
    \correct{$k$}
    \choice{$-k$}
    \choice{$1/k$}
    \choice{$\ln k$}
\end{mcq}
\why{$1/[\ce{A}]_t = kt + 1/[\ce{A}]_0$ --- positive slope for second
order.}

\begin{mcq}
  For which order is the half-life independent of the initial
  concentration?
  \choices
    \choice{zero}
    \correct{first}
    \choice{second}
    \choice{all of them}
\end{mcq}
\why{$t_{1/2} = 0.693/k$ contains no concentration term.}

\begin{mcq}
  A first-order reaction has $t_{1/2} = \SI{25}{\second}$. What fraction of
  the reactant remains after \SI{75}{\second}?
  \choices
    \choice{1/2}
    \choice{1/4}
    \correct{1/8}
    \choice{1/16}
\end{mcq}
\why{Three half-lives.}

\begin{mcq}
  For a second-order reaction, each successive half-life is
  \choices
    \choice{shorter than the last}
    \correct{longer than the last}
    \choice{the same}
    \choice{unpredictable}
\end{mcq}
\why{$t_{1/2} = 1/(k[\ce{A}]_0)$ --- as concentration falls, half-life
grows.}

\begin{mcq}
  For a first-order reaction, which quantity is linear in time?
  \choices
    \choice{$[\ce{A}]$}
    \correct{$\ln[\ce{A}]$}
    \choice{$1/[\ce{A}]$}
    \choice{$[\ce{A}]^2$}
\end{mcq}
\why{Which is why concentrations cannot be linearly interpolated.}

% ---- mechanisms ------------------------------------------------------------
\begin{mcq}
  For the elementary step \ce{2 NO2 -> NO3 + NO}, the rate law is
  \choices
    \choice{$k[\ce{NO2}]$}
    \correct{$k[\ce{NO2}]^2$}
    \choice{$k[\ce{NO3}][\ce{NO}]$}
    \choice{cannot be determined without data}
\end{mcq}
\why{For an elementary step --- and only then --- molecularity gives the
order.}

\begin{mcq}
  A species produced in one step of a mechanism and consumed in a later step
  is called
  \choices
    \choice{a catalyst}
    \correct{an intermediate}
    \choice{a reactant}
    \choice{an activated complex}
\end{mcq}
\why{A catalyst is consumed first and regenerated later --- the reverse
order.}

\begin{mcq}
  A reactant that appears in the balanced equation but not in the rate law
  most likely
  \choices
    \correct{participates only after the rate-determining step}
    \choice{is a catalyst}
    \choice{is not really consumed}
    \choice{indicates an experimental error}
\end{mcq}
\why{Speeding up a step past the bottleneck cannot change the overall rate.}

% ---- collision model and catalysis -----------------------------------------
\begin{mcq}
  For a collision to result in reaction, the molecules must have sufficient
  energy AND
  \choices
    \choice{equal masses}
    \correct{proper orientation}
    \choice{opposite charges}
    \choice{the same velocity}
\end{mcq}
\why{Both conditions filter the enormous number of collisions that occur.}

\begin{mcq}
  A reaction is strongly exothermic yet proceeds imperceptibly slowly at
  room temperature. This is possible because
  \choices
    \correct{$E_a$ and $\Delta H$ are independent quantities}
    \choice{the reaction is not spontaneous}
    \choice{$\Delta H$ must actually be positive}
    \choice{no catalyst is present}
\end{mcq}
\why{A large $E_a$ makes a thermodynamically favourable reaction kinetically
slow.}

\begin{mcq}
  A catalyst increases a reaction rate by
  \choices
    \choice{raising the temperature}
    \choice{increasing $\Delta H$}
    \correct{providing a pathway with lower activation energy}
    \choice{shifting the equilibrium toward products}
\end{mcq}
\why{It changes the mechanism, not the thermodynamics.}
\distractor{D}{a catalyst lowers $E_a$ in both directions equally, so
equilibrium is unchanged}

\examsection{II}{Free Response}{2 questions \textbullet\ 20 minutes}{\calculatorok}

\begin{apcnote}[Kinetics equations]
$\ln[\ce{A}]_t - \ln[\ce{A}]_0 = -kt$ \textbullet\
$\dfrac{1}{[\ce{A}]_t} - \dfrac{1}{[\ce{A}]_0} = kt$ \textbullet\
$t_{1/2} = \dfrac{0.693}{k}$ (first order) \textbullet\
$\ln k = -\dfrac{E_a}{R}\left(\dfrac{1}{T}\right) + \ln A$ \textbullet\
$R = \SI{8.314}{\joule\per\mole\per\kelvin}$
\end{apcnote}

% ---- FRQ 1 -----------------------------------------------------------------
\frq{short}{4}{Zumdahl \S12.3--12.4 \textbullet\ rate law and integrated form}

The decomposition \ce{2 H2O2(aq) -> 2 H2O(l) + O2(g)} is first order in
\ce{H2O2}, with $k = \SI{7.30e-4}{\per\second}$ at \SI{20}{\celsius}. A
solution is prepared with $[\ce{H2O2}]_0 = \SI{0.500}{\Molar}$.

\begin{parts}
  \item Calculate the half-life. \workspace[1.8cm]
        \keyonly{\begin{apcanswer}$t_{1/2} = 0.693/7.30\times10^{-4} =
        \SI{949}{\second}$\end{apcanswer}}
  \item Calculate $[\ce{H2O2}]$ after \SI{1800}{\second}.
        \workspace[2.2cm]
        \keyonly{\begin{apcanswer}$\ln[\ce{A}] = -(7.30\times10^{-4})(1800)
        + \ln(0.500) = -1.314 - 0.693 = -2.007$;
        $[\ce{A}] = \SI{0.134}{\Molar}$\end{apcanswer}}
  \item A student estimates the answer to (b) as \SI{0.125}{\Molar} by
        noting that 1800 s is close to two half-lives. Comment on this
        estimate. \answerlines[2]{It is a reasonable approximation --- two
        half-lives is 1898 s, slightly longer than 1800 s, so slightly more
        than one quarter should remain. The exact value 0.134 M is indeed a
        little above 0.125 M.}
\end{parts}

\begin{rubric}
  \rpoint{1}{(a) \SI{949}{\second}}
  \rpoint{2}{(b) 1 pt correct use of the first-order integrated law;
    1 pt \SI{0.134}{\Molar}}
  \rpoint{1}{(c) recognizes the estimate is sound and explains the
    direction of the small discrepancy}
\end{rubric}

% ---- FRQ 2 -----------------------------------------------------------------
\frq{short}{4}{Zumdahl \S12.5--12.7 \textbullet\ mechanism and energy}

The overall reaction \ce{2 NO2 + F2 -> 2 NO2F} has the experimental rate law
rate $= k[\ce{NO2}][\ce{F2}]$. The proposed mechanism is
\begin{align*}
  \ce{NO2 + F2 &-> NO2F + F} &&\text{slow}\\
  \ce{NO2 + F &-> NO2F} &&\text{fast}
\end{align*}

\begin{parts}
  \item Show that the mechanism is consistent with both the overall
        equation and the rate law. \answerlines[3]{The steps sum to
        \ce{2 NO2 + F2 -> 2 NO2F} (F cancels). The slow step is bimolecular
        in \ce{NO2} and \ce{F2}, predicting rate $=
        k[\ce{NO2}][\ce{F2}]$ --- matching experiment.}
  \item Identify the intermediate and explain how you know.
        \answerlines[2]{F --- it is produced in the first step and consumed
        in the second, so it does not appear in the overall equation.}
  \item Explain why the rate law contains $[\ce{NO2}]$ to the first power
        even though the balanced equation has a coefficient of 2.
        \answerlines[3]{Only one \ce{NO2} is involved in the
        rate-determining step. The second \ce{NO2} reacts in the fast step,
        which occurs after the bottleneck and therefore does not influence
        the overall rate.}
\end{parts}

\begin{rubric}
  \rpoint{1}{(a) steps sum correctly, F shown cancelling}
  \rpoint{1}{(a) slow-step rate law derived and matched to experiment}
  \rpoint{1}{(b) identifies F with produced-then-consumed reasoning}
  \rpoint{1}{(c) explains that only the rate-determining step sets the
    order --- ``coefficients don't matter'' alone earns 0}
\end{rubric}

\examtotal

\end{document}
